\section*{Additional Calibration Methods for the Emoji Hunter Task}

\subsection*{Problem \& Approach}
Neural networks are often overconfident---assigning high probability even when wrong. In HW\,2, we trained an MLP classifier on frozen DINOv2 features for binary emoji detection (1200 train / 400 test, 128$\times$128 images). The base model achieved 78.0\% accuracy with ECE = 0.1278. The main assignment explored temperature scaling, MC dropout, and label smoothing. Here we evaluate three additional methods:

\begin{itemize}
    \item \textbf{Platt Scaling} (post-hoc): Fits a logistic regression on the model's raw logits using a held-out validation set, learning both a slope and bias---more flexible than temperature scaling's single parameter. We split training data 80/20, retrained a fresh model, and fit \texttt{sklearn.LogisticRegression} on validation logits.
    \item \textbf{Weight Decay} (train-time): Adds an L2 penalty $\lambda \|w\|^2$ to the loss, encouraging smaller weights. We tested $\lambda \in \{10^{-4}, 10^{-3}, 10^{-2}\}$ with Adam and selected the best ECE.
    \item \textbf{Focal Loss} (train-time): Modifies cross-entropy with a modulating factor $\text{FL}(p_t) = -(1 - p_t)^\gamma \log(p_t)$ that downweights easy examples \cite{lin2017focal}. We tested $\gamma \in \{1.0, 2.0, 3.0\}$.
\end{itemize}

\subsection*{Results}

\begin{table}[h]
\centering
\begin{tabular}{l l c c}
\toprule
\textbf{Method} & \textbf{Type} & \textbf{Accuracy} & \textbf{ECE $\downarrow$} \\
\midrule
Base Model & --- & 78.00\% & 0.1278 \\
\midrule
\multicolumn{4}{l}{\textit{Main Assignment Methods}} \\
\quad Temp Scaling ($T\!=\!3.0$) & Post-hoc & 78.00\% & 0.0373 \\
\quad MC Dropout (30 passes) & Inference & 77.75\% & 0.1022 \\
\quad Label Smoothing ($\alpha\!=\!0.1$) & Train-time & 78.00\% & 0.0397 \\
\midrule
\multicolumn{4}{l}{\textit{Extra Credit Methods}} \\
\quad Platt Scaling & Post-hoc & 77.25\% & 0.0525 \\
\quad Weight Decay ($\lambda\!=\!10^{-2}$) & Train-time & 77.75\% & 0.1066 \\
\quad Focal Loss ($\gamma\!=\!2.0$) & Train-time & 78.75\% & \textbf{0.0363} \\
\bottomrule
\end{tabular}
\caption{Calibration comparison across all methods. ECE = Expected Calibration Error.}
\label{tab:calibration_results}
\end{table}

\subsection*{Discussion}
\textbf{Focal Loss} ($\gamma\!=\!2.0$) achieved the best ECE overall (0.0363), slightly outperforming temperature scaling (0.0373). By downweighting easy examples during training, it inherently discourages overconfident predictions without any post-hoc adjustment.

\textbf{Platt Scaling} (ECE = 0.0525) improved over the base model but underperformed simpler temperature scaling despite having more parameters. The small validation set (240 samples) likely limited its effectiveness.

\textbf{Weight Decay} (ECE = 0.1066) gave only modest improvement. Generic L2 regularization does not directly target calibration, making it the weakest method tested.

\textbf{Takeaway:} Methods that directly penalize overconfidence (focal loss, label smoothing) or post-hoc rescale probabilities (temperature/Platt scaling) substantially outperform generic regularization (weight decay, MC dropout) for calibration.

\begin{thebibliography}{9}
\bibitem{lin2017focal}
T.-Y.\ Lin, P.\ Goyal, R.\ Girshick, K.\ He, and P.\ Doll\'{a}r, ``Focal Loss for Dense Object Detection,'' \textit{arXiv:1708.02002}, 2017.
\end{thebibliography}
